\section{Conclusion}

This paper investigated the generalization of refusal behavior in safety-aligned Large Language Models (LLMs). Our experiments with \qwen confirm that refusal is a highly contagious and generalizable behavior. Fine-tuning on as few as 20 random benign prompts where the model is taught to refuse led to a broad and disproportionate increase in over-refusal across unrelated tasks, particularly for borderline and ambiguous cases. We identified a phase transition in model behavior, moving from a "Sensitized" state with increased over-refusal to a "Collapsed" state where helpfulness is completely overridden. Our qualitative analysis shows that this induction reinforces the model's sensitivity to semantic triggers that are near the safety boundary.

Our work highlights the sensitivity and potential fragility of the safety boundary in current aligned models. While this sensitivity allows for effective safety training, it also makes models prone to broad over-refusal from minimal "incorrect" data. These findings point towards the need for more precise and context-aware alignment techniques that can maintain robust safety without sacrificing helpfulness. Future work should explore strategies for mitigating this "refusal contagion" and developing more nuanced safety boundaries that are less sensitive to broad semantic triggers.
